\begin{tikzpicture}[
        xscale=1.7, % 修改这里控制横向（代数 n）的间距
        yscale=1.1, % 修改这里控制纵向（空间位置 V）的伸展程度
        >=Stealth,
        % 密集随机游走样式
        rw-path/.style={thick, blue!70!black},
        % 动态屏障样式
        barrier-line/.style={thick, red!70!black, dashed},
        % 参考线样式
        ref-line/.style={thin, gray!60, dashed}
    ]

        % 0. 绘制填充阴影 (放在最底层)
        \fill[blue!5] (0,0) 
            -- (0.2, 0.6) -- (0.4, -0.3) -- (0.6, 0.2) -- (0.8, 1.1) 
            -- (1.1, 0.3) -- (1.3, 0.7) -- (1.6, -0.5) -- (1.9, 0.2) 
            -- (2.2, 1.5) -- (2.5, 0.4) -- (2.8, 0.9) -- (3.1, -0.45) 
            -- (3.4, 0.1) -- (3.7, -0.4) -- (4.0, 0.6) -- (4.3, -0.2) -- (4.5, 0.3) % 蓝色路径结束
            -- (4.5, -0.5) -- (2.2, -0.5) -- (2.1, -0.9) -- (0.8, -0.9) -- (0.7, -1.4) -- (0.2, -1.4) -- (0, -2.0) % 红色路径逆序
            -- cycle;

        % 1. 坐标轴
        \draw[->, thick] (0, -2.5) -- (0, 2.2) node[above] {$Y_n$};
        \draw[->, thick] (0, 0) -- (5.5, 0) node[right] {$n$};

        % 2. 终点限制线 y = Delta/c (假设 Delta/c = 0.8)
        \draw[thick, black!80] (5, 0.8) -- (5.4, 0.8) node[right] {$\frac{\Delta}{c}$};
        \draw[ref-line] (0, 0.8) -- (5, 0.8);

        % 3. 时间 L 的垂线
        \draw[dashed, black!80] (4.5, -2.5) -- (4.5, 2.0) node[above] {$L \ge \Delta^{2+c}$};

        % 4. 密集随机游走路径 (零均值特征)
        % 定义一组波动点
        \draw[rw-path] (0,0) 
            -- (0.2, 0.6) -- (0.4, -0.3) -- (0.6, 0.2) -- (0.8, 1.1) 
            -- (1.1, 0.3) -- (1.3, 0.7) -- (1.6, -0.5) -- (1.9, 0.2) 
            -- (2.2, 1.5) -- (2.5, 0.4) -- (2.8, 0.9) -- (3.1, -0.45) 
            -- (3.4, 0.1) -- (3.7, -0.4) -- (4.0, 0.6) -- (4.3, -0.2) -- (4.5, 0.3);

        % 5. 动态屏障 Max Y_i - Delta (假设 Delta = 2.0)
        % 初始 Max=0 -> 屏障 -2.0
        % 第一高峰 0.6 -> 屏障 -1.4
        % 第二高峰 1.1 -> 屏障 -0.9
        % 第三高峰 1.5 (最高) -> 屏障 -0.5
        \draw[barrier-line] 
            (0, -2.0) -- (0.2, -1.4) -- (0.7, -1.4) -- (0.8, -0.9) 
            -- (2.1, -0.9) -- (2.2, -0.5) -- (4.5, -0.5) 
            node[right] {$\max_{i\le n} Y_i - \Delta$};

        % 6. 标注 Delta 间距 (在最高点处)
        \draw[ref-line] (0, 1.5) -- (2.2, 1.5);
        \draw[<->] (2.3, 1.5) -- (2.3, -0.5) node[midway, left] {$\Delta$};

        % 7. 关键点标注
        \node[below left] at (0,0) {$0$};
        \fill[blue!70!black] (4.5, 0.3) circle (1.5pt) node[right, black] {$Y_L$};

    \end{tikzpicture}